% !PRINTED 183
\weberpage{183}

Die Darstellung ist, wie (2) zeigt, auf unendlich viele verschiedene Arten möglich, da die ganz willkürlichen Grössen $x'$ noch in dieser Formel vorkommen.

Eine Darstellung durch weniger als $m$ Quadrate ist nur dann möglich, wenn die Determinante der Function $\varphi$ verschwindet.

Da in dem Ausdruck (2) für $Y_1$ eine Quadratwurzel vorkommt, so kann $Y_1$ reell oder imaginär sein, selbst wenn die Co\"efficienten von $\varphi$ und die $x$ und $x'$ reell vorausgesetzt werden.

Setzen wir aber in dem Falle, wo $Y_1$ imaginär ist, $iY_1$ an Stelle von $Y_1$, so können wir unseren Satz auch so aussprechen:

\gesperrt{Eine reelle quadratische Form von $m$ Veränderlichen lässt sich als Summe von höchstens $m$ positiven oder negativen Quadraten reeller linearer Functionen der $x$ darstellen}\footnote{Für die allgemeine Theorie der linearen Transformation der Formen zweiten Grades sind ausser den Untersuchungen von \gesperrt{Jacobi} und \gesperrt{Hesse} (\gesperrt{Hesse}, „Vorlesungen über analytische Geometrie des Raumes“, dritte Auflage, besorgt von \gesperrt{Gundelfinger}, Leipzig 1876) besonders die Arbeiten von \gesperrt{Weierstrass} und \gesperrt{Kronecker} (Monatsberichte der Berliner Akademie 1868, 1874) von Bedeutung. Eingehende historische Darlegung der Entwickelung der Frage in dem „Bericht über den gegenwärtigen Stand der Invariantentheorie“ von \gesperrt{Franz Meyer} (erster Jahresbericht der deutschen Mathematiker-Vereinigung, Berlin 1892).}.

\begin{center}
\S. 58.

\gesperrt{Trägheitsgesetz der quadratischen Formen.}
\end{center}

Für die Zerlegung einer quadratischen Form in Quadrate gilt nun der folgende Satz, der von \gesperrt{Sylvester} den Namen des Gesetzes der Trägheit der quadratischen Formen erhalten hat (Philos. Mag. 1852).

\gesperrt{Wie man auch eine reelle quadratische Form $\varphi(x)$ in die Summe von positiven und negativen Quadraten linearer Functionen zerlegen mag, die Anzahl der positiven und negativen dieser Quadrate, und also auch ihre Gesammtzahl, ist immer dieselbe, vorausgesetzt, dass zwischen diesen linearen Functionen keine lineare Abhängigkeit besteht.}

Der Beweis dieses wichtigen Satzes ist sehr einfach.

Es seien
%CONT>>
